\documentclass[17pt]{article}

\usepackage{cmap}	% поиск по документу
\usepackage{hyperref}	% поддержка гиперссылок
\usepackage{geometry}	% настройка геометрии документа
\usepackage{multicol}   % настройка шаблона расположения текста
\usepackage{extsizes}	% расширенная настройка размеров шрифтов
\usepackage{indentfirst}    % отступ для первого абзаца
\usepackage{tabularx}	% расширенные таблицы
\usepackage{makecell}   % настройка клеток таблиц
\usepackage{caption}    % настройка подписей к вставкам
\usepackage[table]{xcolor}  % покраска таблиц
\usepackage{graphicx}	% вставка изображений
\usepackage{wrapfig, booktabs}	% оформление таблиц
\usepackage{graphicx, caption, subcaption}
\usepackage{minted}    % вставка кода

\usepackage{color}	% настройка цвета текста - \textcolor
\usepackage{soul}	% выделение текста - \hl
\usepackage[T2A]{fontenc}	% поддержка кириллических символов
\usepackage[utf8]{inputenc} % поддержка кириллических символов
\usepackage[english,russian]{babel}	% поддержка кириллических символов
\usepackage{csquotes, biblatex}   % вывод библиографии

\usepackage{mathcmd}	% математические команды
\usepackage{amsfonts}	% математические символы 1
\usepackage{amssymb}	% математические символы 2
\usepackage{amsmath}	% математические символы 3
\usepackage{amsthm}	% работа с теоремами

% Настройка окружения теорем 1 (RU)
\newtheorem{theorem}{Теорема}[section]
\newtheorem{corollary}{Следствие}[theorem]
\newtheorem{lemma}[theorem]{Лемма}
\newtheorem{example}{Пример}[section]
\newtheorem*{remark}{Утверждение}
\addto\captionsrussian{\renewcommand\proofname{Док-во:}}

% Настройка окружения теорем 2 (EN)
%\newtheorem{theorem}{TH}[section]
%\newtheorem{corollary}{CO}[theorem]
%\newtheorem{lemma}[theorem]{LM}
%\newtheorem{example}{EX}[section]
%\addto\captionsrussian{\renewcommand\proofname{PF:}}

\renewcommand{\labelitemii}{$\circ$}
\renewcommand{\labelitemiii}{-}

\addbibresource{biblio.bib}
\graphicspath{ {./images/part1/} }

% Настройка оформления вставок кода
\setminted{xleftmargin=\parindent, fontsize=\footnotesize, linenos, tabsize=2, breaklines}
\usemintedstyle{emacs} 

% Настройка отступов
\geometry{
    a4paper,
    % total={200mm,287mm},
    left=10mm,
    right=10mm,
    top=10mm,
    bottom=20mm
}

% Заголовок
\title{
    Методы моделирования и анализа стохастических систем \\
    \large Лабораторная работа № 3, Часть 1 \\
    Анализ детерминированной задачи
}
\author{Костарев Иван (МЕНМ-150201)}
\date{весна 2025/2026}

\begin{document}

\maketitle
\tableofcontents
\newpage



% \section*{Анализ детерминированной задачи}
% Определим прогроаммный класс двумерной динамической системы \texttt{DynamicSystem2D} с возможностью нахождения решения методом Рунге-Кутты 4 порядка с суммарной ошибкой $O(h^4)$ на интервале решения, с возможностью построения бифуркационной диаграммы и фазового портрета системы, нахождения равновесий и предельных циклов.

\section{Нейронная модель ФитцХью-Нагумо (FHN)}

$$
\begin{cases}
    \dot{x} = \frac{1}{\delta} (x - \frac{x^3}{3} - y) \\
    \dot{y} = x + a
\end{cases}
$$

1. Найдем равновесия системы.

$$
\begin{cases}
    \frac{1}{\delta} (x - \frac{x^3}{3} - y) = 0 \\
    x + a = 0
\end{cases} \implies
\begin{cases}
    \overline{x} = -a \\
    \overline{y} = \frac{a^3}{3} - a
\end{cases}
$$

2. Исследуем равновесия на устойчивость.

$$
J(x, y) = \begin{pmatrix}
    \frac{1}{\delta} (1 - x^2) & \frac{1}{\delta} \\
    1 & 0
\end{pmatrix}
$$

$$
|J(x, y) - \lambda E| = 0
$$

$$
\lambda_{1,2}(x, y) = \frac{-\delta x^2 + \delta \pm \delta \sqrt{x^4 - 2 x^2 - 4 \delta + 1}}{2 \delta^2}
$$

Далее с помощью численной процедуры исследуем поведение собственных чисел якобиана в равновесных точках при различных значениях параметра $a$.

% \begin{minted}{python}
% FHN = DynamicSystem2D(
%     model_FHN,
%     params=dict(delta=0.1, a=1)
% )
% \end{minted}

3. Построим бифуркационную диаграмму по параметру $a$ (Рис. \ref{fig:fhn_bifurcation_diagram}).

\begin{figure}[H]
    \centering
    \includegraphics[width=0.6\textwidth]{fhn_bifurcation_diagram.png}
    \caption{Бифуркационная диаграмма системы FHN.}
    \label{fig:fhn_bifurcation_diagram}
\end{figure}


4. Построим типичные фазовые портреты для различных значений $a$ по бифуркационной диаграмме (Рис \ref{fig:fhn_phase_portraits}).

\begin{figure}[H]
    \centering
    \includegraphics[width=1.0\textwidth]{fhn_phase_portraits.png}
    \caption{Фазовые портреты системы FHN.}
    \label{fig:fhn_phase_portraits}
\end{figure}

Для нейронной модели ФитцХью-Нагумо при значении параметра $a = 1$ наблюдается бифуркация Хопфа.
При $a < 1$ в системе существует устойчивый предельный цикл, и равновесие $(\overline{x}, \overline{y})$ неустойчиво.
При $a > 1$ предельный цикл исчезает, и равновесие $(\overline{x}, \overline{y})$ становится устойчивым.



\newpage
\section{Климатическая модель Зальцман-Николис (SN)}

$$
\begin{cases}
    \dot{x} = y - x \\
    \dot{y} = -a x + b y - x^2 y
\end{cases}
$$

1. Найдем равновесия системы.

\begin{multline*}
    \begin{cases}
        y - x = 0 \\
        -a x + b y - x^2 y = 0
    \end{cases}
    \implies
    \begin{cases}
        x = y \\
        x(x^2 + a - b) = 0
    \end{cases}
    \implies \\
    \implies
    \begin{cases}
        (\overline{x}_1, \overline{y}_1) = (0, 0) \\
        (\overline{x}_2, \overline{y}_2) = (\sqrt{b - a}, \sqrt{b - a}) \\
        (\overline{x}_3, \overline{y}_3) = (-\sqrt{b - a}, -\sqrt{b - a})
    \end{cases}
\end{multline*}

2. Исследуем равновесия на устойчивость.

$$
J(x, y) = \begin{pmatrix}
    -1 & 1 \\
    -a - 2xy & b - x^2
\end{pmatrix}
$$

$$
|J(x, y) - \lambda E| = 0
$$

$$
\lambda_{1,2}(x, y) = \frac{\pm \sqrt{-4a + b^2 - 2b x^2 + 2b + x^4 - 2x^2 - 8xy + 1} + b - x^2 - 1}{2}
$$

Далее с помощью численной процедуры исследуем поведение собственных чисел якобиана в равновесных точках при различных значениях параметра $a$.

% \begin{minted}{python}
% SN = DynamicSystem2D(
%     model_SN,
%     params=dict(a=1, b=2)
% )
% \end{minted}

3. Построим бифуркационную диаграмму по параметру $a$ (Рис. \ref{fig:sn_bifurcation_diagram}). Изобразим предельные циклы и равновесия:

$$
\begin{cases}
    (\overline{x}_1, \overline{y}_1) = (0, 0) \\
    (\overline{x}_2, \overline{y}_2) = (\sqrt{b - a}, \sqrt{b - a}) \\
    (\overline{x}_3, \overline{y}_3) = (-\sqrt{b - a}, -\sqrt{b - a})
\end{cases}
$$

\begin{figure}[H]
    \centering
    \includegraphics[width=0.6\textwidth]{sn_bifurcation_diagram_single.png}
    \caption{Бифуркационная диаграмма системы SN.}
    \label{fig:sn_bifurcation_diagram}
\end{figure}


4. Построим типичные фазовые портреты для различных значений $a$ по бифуркационной диаграмме (Рис. \ref{fig:sn_phase_portraits}).

\begin{figure}[H]
    \centering
    \includegraphics[width=1.0\textwidth]{sn_phase_portraits.png}
    \caption{Фазовые портреты системы SN.}
    \label{fig:sn_phase_portraits}
\end{figure}

Для климатической модели Зальцман-Николис при значении параметра $a = 1$ наблюдается бифуркация Хопфа: равновесия $(\overline{x}_2, \overline{y}_2)$, $(\overline{x}_3, \overline{y}_3)$ становятся неустойчивыми.
При $a < 1$ в системе существуют неустойчивые предельные циклы вокруг устойчивых фокусов $(\overline{x}_2, \overline{y}_2)$, $(\overline{x}_3, \overline{y}_3)$.
При $a > 1$ фокусы становятся неустойчивыми, и предельные циклы фокруг них исчезают.
При $a = 2$ наблюдается бифуркация типа ``вилка'': происходит слияние седла $(\overline{x}_1, \overline{y}_1)$ с неустойчивыми фокусами $(\overline{x}_2, \overline{y}_2)$, $(\overline{x}_3, \overline{y}_3)$.
При всех рассмотренных значениях параметра $a$ в системе сохраняется внешний устойчивый предельный цикл.

\end{document}
